\documentclass[10pt,a4paper]{amsart}
\usepackage{amsmath,amssymb,amsthm,graphicx}
% Retained minimal preamble for the Stacks Project chapter more-algebra.tex (More on Algebra).
% The chapter file pinned-original-source.tex carries only \input{preamble} on its first line, so no
% preamble text lives inside the chapter itself.  The Stacks Project uses one single shared preamble for
% every chapter; the snapshot stages it identically as sources/stacks-chow/preamble.tex and
% sources/stacks-etale-cohomology/preamble.tex.  The block below is copied mechanically, line for line,
% from sources/stacks-chow/preamble.tex lines 14-18 (the Xy-pic setup) and 154-217 (the native theorem
% declarations together with the mathematical macros).  Nothing was added, removed or redefined; the
% wrapper class and the amsmath/amssymb/amsthm/graphicx base packages are supplied by the assembler.
% Chapter: more-algebra.tex -- More on Algebra.
% For commutative diagrams we use Xy-pic
\usepackage[all]{xy}

% We use 2cell for 2-commutative diagrams.
\xyoption{2cell}

% Theorem environments.
%
\theoremstyle{plain}
\newtheorem{theorem}[subsection]{Theorem}
\newtheorem{proposition}[subsection]{Proposition}
\newtheorem{lemma}[subsection]{Lemma}

\theoremstyle{definition}
\newtheorem{definition}[subsection]{Definition}
\newtheorem{example}[subsection]{Example}
\newtheorem{exercise}[subsection]{Exercise}
\newtheorem{situation}[subsection]{Situation}

\theoremstyle{remark}
\newtheorem{remark}[subsection]{Remark}
\newtheorem{remarks}[subsection]{Remarks}

\numberwithin{equation}{subsection}

% Macros
%
\def\lim{\mathop{\mathrm{lim}}\nolimits}
\def\colim{\mathop{\mathrm{colim}}\nolimits}
\def\Spec{\mathop{\mathrm{Spec}}}
\def\Hom{\mathop{\mathrm{Hom}}\nolimits}
\def\Ext{\mathop{\mathrm{Ext}}\nolimits}
\def\SheafHom{\mathop{\mathcal{H}\!\mathit{om}}\nolimits}
\def\SheafExt{\mathop{\mathcal{E}\!\mathit{xt}}\nolimits}
\def\Sch{\mathit{Sch}}
\def\Mor{\mathop{\mathrm{Mor}}\nolimits}
\def\Ob{\mathop{\mathrm{Ob}}\nolimits}
\def\Sh{\mathop{\mathit{Sh}}\nolimits}
\def\NL{\mathop{N\!L}\nolimits}
\def\CH{\mathop{\mathrm{CH}}\nolimits}
\def\proetale{{pro\text{-}\acute{e}tale}}
\def\etale{{\acute{e}tale}}
\def\QCoh{\mathit{QCoh}}
\def\Ker{\mathop{\mathrm{Ker}}}
\def\Im{\mathop{\mathrm{Im}}}
\def\Coker{\mathop{\mathrm{Coker}}}
\def\Coim{\mathop{\mathrm{Coim}}}

% Boxtimes
%
\DeclareMathSymbol{\boxtimes}{\mathbin}{AMSa}{"02}

%
% Macros for moduli stacks/spaces
%
\def\QCohstack{\mathcal{QC}\!\mathit{oh}}
\def\Cohstack{\mathcal{C}\!\mathit{oh}}
\def\Spacesstack{\mathcal{S}\!\mathit{paces}}
\def\Quotfunctor{\mathrm{Quot}}
\def\Hilbfunctor{\mathrm{Hilb}}
\def\Curvesstack{\mathcal{C}\!\mathit{urves}}
\def\Polarizedstack{\mathcal{P}\!\mathit{olarized}}
\def\Complexesstack{\mathcal{C}\!\mathit{omplexes}}
% \Pic is the operator that assigns to X its picard group, usage \Pic(X)
% \Picardstack_{X/B} denotes the Picard stack of X over B
% \Picardfunctor_{X/B} denotes the Picard functor of X over B
\def\Pic{\mathop{\mathrm{Pic}}\nolimits}
\def\Picardstack{\mathcal{P}\!\mathit{ic}}
\def\Picardfunctor{\mathrm{Pic}}
\def\Deformationcategory{\mathcal{D}\!\mathit{ef}}
\usepackage{cleveref}
\crefname{originallemma}{original lemma}{original lemmas}
\Crefname{originallemma}{Original lemma}{Original lemmas}
\makeatletter
\newlabel{fields-lemma-Kummer}{{09DX}{source}{}{sourcepointer-1}{}}
\newlabel{fields-lemma-Kummer@cref}{{[originallemma][0][]09DX}{[1][0][]source}{}{}{}}
\makeatother
\crefname{originallemma}{original lemma}{original lemmas}
\Crefname{originallemma}{Original lemma}{Original lemmas}
\makeatletter
\newlabel{lemma-make-degree-q-extension}{{09EW}{source}{}{sourcepointer-2}{}}
\newlabel{lemma-make-degree-q-extension@cref}{{[originallemma][0][]09EW}{[1][0][]source}{}{}{}}
\makeatother
\crefname{originallemma}{original lemma}{original lemmas}
\Crefname{originallemma}{Original lemma}{Original lemmas}
\makeatletter
\newlabel{lemma-weakly-unramified-goes-up-along-totally-ramified}{{09ER}{source}{}{sourcepointer-3}{}}
\newlabel{lemma-weakly-unramified-goes-up-along-totally-ramified@cref}{{[originallemma][0][]09ER}{[1][0][]source}{}{}{}}
\makeatother
\crefname{originallemma}{original lemma}{original lemmas}
\Crefname{originallemma}{Original lemma}{Original lemmas}
\makeatletter
\newlabel{lemma-pre-purely-inseparable-case}{{09EX}{source}{}{sourcepointer-4}{}}
\newlabel{lemma-pre-purely-inseparable-case@cref}{{[originallemma][0][]09EX}{[1][0][]source}{}{}{}}
\makeatother
\title{Producing a totally ramified finite Galois extension that either solves the ramification weakly or puts the degree p extension on finite level}
\author{The Stacks Project Authors}
\date{stacks-project:more-algebra.tex:local-snapshot-20261012; GNU FDL 1.2 or later}
\begin{document}
\maketitle
\setcounter{subsection}{1}
\noindent\textit{Editorial note.} All delivered mathematical text is the author's own, copied from the pinned Stacks Project chapter file more-algebra.tex; no statement or proof step was rewritten, reordered or removed. The counted statement is the labeled lemma environment at source lines 35087-35100 (canonical tag 09F4) and the counted proof is the proof environment at source lines 35102-35140, which immediately follows it and is its own proof.\par\medskip
\noindent\textit{Editorial note.} The proof refers four times, three times to earlier lemmas of this chapter and once to Kummer theory in another chapter; all four are navigation stubs with their roles recorded in the reasons.\par\medskip
\noindent\textit{Editorial note.} The statement and prose use the terminology of weak solutions and of extensions of finite level defined earlier in the same section; those definitions lie outside the counted spans, and the disjunctive conclusion is stated in that language exactly as the author wrote it.\par\medskip
\noindent\textit{Editorial note.} The chapter file pinned-original-source.tex carries only an input of the shared preamble on its first line, so no preamble text lives inside the chapter. The retained file sources/stacks-more-algebra/retained-preamble.tex is a mechanical, line-for-line copy of the project's single shared preamble as staged in the snapshot (Xy-pic setup plus the native theorem declarations and mathematical macros). Nothing was added, removed or redefined, and the excerpt serves only as preamble, never as a span coordinate.\par\medskip
\noindent\textit{Editorial note.} The source URL is the source-level Stacks Project URL, as required when more than one card is produced from one source; the per-claim tag is carried by claim\_key and by this note. The counted statement and proof cite no bibliography entry, so this card needs no bibliography.\par\medskip
\noindent\textit{Allowed standard prerequisites (editorial).}\begin{itemize}
\item Kummer theory: a degree p Galois extension of a field containing a primitive p-th root of unity is obtained by adjoining the p-th root of an element of the base field. invoked at proof lines 35103-35105. Produces the element b whose p-th root generates M over L, which is the input to every case distinction in the proof.
\item Weakly unramified extensions go up along a totally ramified extension, so the integral closure in the compositum remains weakly unramified over the new base. invoked at proof lines 35113-35115. Shows that the compositum of the chosen totally ramified extension with the given towers again has the weak unramifiedness required of the conclusion.
\item The pre-purely-inseparable case: if the residue class of the unit is not a p-th power then the extension of discrete valuation rings is weakly unramified. invoked at proof lines 35128-35130. Disposes of the case in which the residue class of the unit is not a p-th power, which is one of the two alternatives in the conclusion.
\end{itemize}\medskip
\ifdefined\hypertarget\hypertarget{sourcepointer-1}{}\fi
\noindent\textit{Source reference.} Original lemma 09DX: sources/stacks-chow/tags.txt, source lines 11880--11880. The reference is Kummer theory, invoked at proof lines 35103-35105 to write the top extension as adjoining a p-th root of an element of L; the whole case analysis built on that element is inside the counted span.\par\medskip
\ifdefined\hypertarget\hypertarget{sourcepointer-2}{}\fi
\noindent\textit{Source reference.} Original lemma 09EW: sources/stacks-more-algebra/pinned-original-source.tex, source lines 34641--34641. The reference supplies the existence of the degree p separable totally ramified extension used at proof lines 35108-35111; the case analysis that consumes it is inside the counted span.\par\medskip
\ifdefined\hypertarget\hypertarget{sourcepointer-3}{}\fi
\noindent\textit{Source reference.} Original lemma 09ER: sources/stacks-more-algebra/pinned-original-source.tex, source lines 34490--34490. The reference is applied at proof lines 35113-35115 to see that the integral closure in the compositum is weakly unramified over the new base; the resulting alternative of the conclusion is produced inside the counted span.\par\medskip
\ifdefined\hypertarget\hypertarget{sourcepointer-4}{}\fi
\noindent\textit{Source reference.} Original lemma 09EX: sources/stacks-more-algebra/pinned-original-source.tex, source lines 34663--34663. The reference is applied at proof lines 35128-35130 to conclude weak unramifiedness when the residue class of the unit is not a p-th power; the complementary case, which yields finite level, is handled inside the counted span.\par\medskip
\begin{lemma}
\label{lemma-make-finite-level}
Let $A \subset B \subset C$ be extensions of discrete valuation rings
with fractions fields $K \subset L \subset M$. Assume that
\begin{enumerate}
\item $A$ has mixed characteristic $(0, p)$,
\item $A \subset B$ is weakly unramified,
\item $B$ contains a primitive $p$th root of $1$, and
\item $M/L$ is Galois of degree $p$.
\end{enumerate}
Then there exists a finite Galois extension $K_1/K$ totally ramified
with respect to $A$ which is either a weak solution for $A \to C$
or is such that $M_1/L_1$ is a degree $p$ extension of finite level.
\end{lemma}

\begin{proof}
Let $\pi \in A$ be a uniformizer. By Kummer theory
(Fields, Lemma \ref{fields-lemma-Kummer}) $M$ is obtained
from $L$ by adjoining the root of $y^p = b$ for some $b \in L$.

\medskip\noindent
If $\text{ord}_B(b)$ is prime to $p$, then we choose a degree $p$
separable extension $K_1/K$
totally ramified with respect to $A$ (for example
using Lemma \ref{lemma-make-degree-q-extension}).
Let $A_1$ be the integral closure of $A$ in $K_1$.
By Lemma \ref{lemma-weakly-unramified-goes-up-along-totally-ramified}
the integral closure $B_1$ of $B$ in $L_1 = L \otimes_K K_1$
is a discrete valuation ring weakly unramified over $A_1$.
If $K_1/K$ is not a weak solution for $A \to C$, then
the integral closure $C_1$ of $C$ in $M_1 = M \otimes_K K_1$ is a
discrete valuation ring and $B_1 \to C_1$ has ramification index $p$.
In this case, the field $M_1$ is obtained from $L_1$ by adjoining the
$p$th root of $b$ with $\text{ord}_{B_1}(b)$ divisible by $p$.
Replacing $A$ by $A_1$, etc we may assume that $b = \pi^n u$ where
$u \in B$ is a unit and $n$ is divisible by $p$. Of course, in
this case the extension $M$ is obtained from $L$ by adjoining
the $p$th root of a unit.

\medskip\noindent
Suppose $M$ is obtained from $L$ by adjoining the root of
$y^p = u$ for some unit $u$ of $B$. If the residue class of $u$
in $\kappa_B$ is not a $p$th power, then $B \subset C$ is
weakly unramified (Lemma \ref{lemma-pre-purely-inseparable-case})
and we are done. Otherwise, we can replace our choice of $y$ by
$y/v$ where $v^p$ and $u$ have the same image in $\kappa_B$.
After such a replacement we have
$$
y^p = 1 + \pi b
$$
for some $b \in B$. Then we see that $P(z) = \pi b/ w^p$ where
$z = (y - 1)/w$. Thus we see that the extension is a degree $p$
extension of finite level with $\xi = \pi b / w^p$.
\end{proof}

\end{document}
